\chapter{Podsumowanie}
\label{ch:podsumowanie}

Rozdział zamyka pracę. W~podrozdziale \ref{sec:podsumowanie_praca} zestawiono
wykonane zadania z~celem postawionym w~rozdziale~\ref{ch:wstep}, oceniono stopień
jego realizacji i~sformułowano najważniejsze wnioski z~analizy porównawczej.
Podrozdział \ref{sec:podsumowanie_rozwoj} wskazuje możliwości dalszego rozwoju
--- zarówno samych algorytmów, jak i~środowiska symulacyjnego oraz metodyki
badań --- a~podrozdział \ref{sec:podsumowanie_zastosowania} omawia potencjalne
obszary zastosowania uzyskanych wyników i~opracowanego oprogramowania.

\section{Podsumowanie pracy}
\label{sec:podsumowanie_praca}

Celem pracy było opracowanie, implementacja oraz przeprowadzenie statystycznej
analizy porównawczej wielokryterialnych algorytmów metaheurystycznych NSGA-II,
NSGA-III i~GWO w~zadaniu planowania tras zespołu bezzałogowych statków
powietrznych w~niestacjonarnym środowisku. Cel ten zrealizowano w~następujących
krokach.

Zaprojektowano i~zaimplementowano autorskie środowisko symulacyjne
(rozdział~\ref{ch:wyniki}): dyskretną siatkę trójwymiarową z~proceduralnie
generowaną rzeźbą terenu, strefami zakazu lotu, modelem wykrywalności radarowej
oraz zmiennym w~czasie polem wiatru. Trasy zespołu dronów oceniane są według
trzech minimalizowanych kryteriów --- czasu wykonania misji, zużycia energii
i~ryzyka wykrycia --- przy twardym egzekwowaniu ograniczeń przestrzennych
i~kolizyjnych, bez funkcji kary zniekształcających wartości kryteriów.

Zaimplementowano trzy algorytmy na wspólnym kodowaniu rozwiązań, ze wspólnym
ewaluatorem i~jednakowym budżetem ocen funkcji celu (rozdział~\ref{ch:sota}),
co wyeliminowało różnice implementacyjne jako źródło rozbieżności wyników.
Ponieważ klasyczny GWO jest algorytmem jednokryterialnym operującym w~przestrzeni
ciągłej, opracowano jego wielokryterialny i~dyskretny wariant: MOGWO z~zewnętrznym
archiwum Pareto, ruletkowym wyborem liderów oraz aktualizacją pozycji opartą na
głosowaniu propozycji liderów dla poszczególnych genów trasy.

Zbudowano protokół porównawczy oparty na dokładnie obliczanej hiperobjętości
ze wspólną normalizacją frontów oraz na testach statystycznych dla prób
sparowanych (test Friedmana, testy Wilcoxona z~poprawką Holma, delta Cliffa).
Poprawność metryki zweryfikowano dwiema niezależnymi implementacjami, a~cały
eksperyment --- 20 instancji problemu w~eksperymencie głównym oraz seria
czterech budżetów obliczeniowych --- jest w~pełni odtwarzalny dzięki jawnemu
ziarnowaniu generatorów losowych i~skryptom automatyzującym obliczenia
i~analizę.

Cel pracy uznaje się za zrealizowany w~zakresie określonym w~rozdziale
\ref{ch:wyniki}: porównanie przeprowadzono dla planowania z~pełną znajomością
modelu środowiska, w~którym niestacjonarność przejawia się zależnością kosztu
ruchu od czasu jego wykonania. Adaptacja tras w~trakcie lotu do zmian
nieznanych w~chwili planowania nie była przedmiotem pracy i~stanowi jeden
z~kierunków jej rozwoju.

\subsection{Najważniejsze wnioski}
\label{sec:podsumowanie_wnioski}

\begin{enumerate}
    \item Algorytmy ewolucyjne NSGA-II i~NSGA-III uzyskują istotnie wyższą
    hiperobjętość niż dyskretny wariant MOGWO ($p_{\text{Holm}} < 0{,}001$,
    $\delta$ Cliffa $0{,}81$--$0{,}88$); MOGWO przegrywa w~19 z~20 instancji
    i~nie wygrał w~żadnej.
    \item Różnica między NSGA-II (średnia $0{,}73$) a~NSGA-III ($0{,}58$) nie jest
    istotna statystycznie ($p = 0{,}114$, efekt mały) i~zależy od budżetu:
    NSGA-III jest lepszy przy małym budżecie, a~przy dużym oba algorytmy dają
    fronty tej samej jakości. Jest to zgodne z~doniesieniami literaturowymi
    o~braku przewagi NSGA-III dla trzech kryteriów.
    \item Główną słabością MOGWO jest kryterium czasu: w~$30\%$ instancji wataha
    nie znalazła tras doprowadzających wszystkie drony do celu, a~jej krzywa
    zbieżności wypłaszcza się po $60\%$ budżetu, co wskazuje na przedwczesną
    utratę różnorodności.
    \item Duży rozmiar archiwum MOGWO (średnio $56$ rozwiązań) nie przekłada
    się na jakość frontu --- liczność zbioru niezdominowanego nie jest miarą jego
    jakości.
    \item Jakość frontów wszystkich algorytmów rośnie bez nasycenia w~całym
    badanym zakresie budżetu ($100$--$1000$ iteracji); przy równym budżecie ocen
    MOGWO potrzebuje $1{,}8$ raza więcej czasu procesora.
    \item Zaobserwowane różnice dają się wyjaśnić niedopasowaniem genowego
    operatora MOGWO do kodowania z~zależnościami pozycyjnymi i~czasowymi,
    brakiem elitaryzmu w~watasze oraz wygaszaniem eksploracji przez harmonogram
    parametru $a$ --- co wskazuje konkretne kierunki dalszych prac nad adaptacją
    GWO do dyskretnego planowania tras.
\end{enumerate}

Odpowiedź na pytanie o~przydatność badanych metod jest zatem zróżnicowana.
Algorytmy z~rodziny NSGA okazały się skutecznym narzędziem wyznaczania zbioru
tras kompromisowych dla zespołu dronów w~rozważanym środowisku, a~wybór między
NSGA-II i~NSGA-III ma mniejsze znaczenie niż dobór budżetu obliczeniowego.
Bezpośrednie przeniesienie mechanizmów GWO na grunt dyskretny, mimo zachowania
wszystkich jego charakterystycznych elementów, nie dało konkurencyjnej metody
--- wynik ten jest jednak wartościowy poznawczo, gdyż pozwolił zidentyfikować,
które z~założeń algorytmu ciągłego nie przenoszą się na kodowanie trasowe.

\section{Możliwości dalszego rozwoju}
\label{sec:podsumowanie_rozwoj}

\paragraph{Rozwój algorytmów.}
Analiza z~rozdziału~\ref{ch:analiza} wskazuje trzy modyfikacje MOGWO, które
bezpośrednio adresują zidentyfikowane słabości: przejmowanie od liderów spójnych
segmentów tras zamiast pojedynczych genów (zachowanie kontekstu przestrzennego
i~czasowego), wprowadzenie elitaryzmu w~watasze przez zachowanie części archiwum
w~populacji oraz dolne ograniczenie prawdopodobieństwa mutacji niezależne od
parametru $a$. Warto również zweryfikować wpływ zastąpienia odległości zatłoczenia
oryginalną siatką hipersześcianów. Niezależnie od GWO, naturalnym rozszerzeniem
jest włączenie do porównania innych paradygmatów --- dekompozycyjnego MOEA/D,
wielokryterialnej optymalizacji rojem cząstek czy algorytmów mrówkowych,
które ze względu na konstrukcyjny charakter dobrze pasują do kodowania trasowego
--- a~także hybrydyzacja z~lokalnym przeszukiwaniem lub operatorami naprawczymi
wykorzystującymi wiedzę o~problemie (np. krzyżowanie w~punkcie, w~którym trasy
rodziców przechodzą przez tę samą komórkę). Osobnym zadaniem jest systematyczne
strojenie parametrów operatorów, które w~niniejszej pracy przyjęto na podstawie
literatury, bez optymalizacji pod rozważany problem.

\paragraph{Rozwój środowiska symulacyjnego.}
Najpilniejszym uzupełnieniem jest eksperyment ablacyjny z~wyłączonymi podmuchami
wiatru, który pozwoliłby ilościowo wyodrębnić wpływ niestacjonarności na jakość
frontów. Dalsze kroki to wprowadzenie zmian środowiska nieznanych w~chwili
planowania --- dryfu lub pojawiania się radarów, ruchomych przeszkód --- wraz
z~planowaniem w~trybie przesuwanego horyzontu, w~którym algorytm przeplanowuje
trasę co określoną liczbę tików na podstawie bieżącej populacji. Model można
również wzbogacić o~egzekwowanie budżetu energii jako ograniczenia zasięgu,
zróżnicowane prędkości dronów, ciągłą (pozbawioną składnika stałego) funkcję
kosztu ryzyka radarowego oraz większe zespoły i~mapy w~celu zbadania
skalowalności. Ponieważ ewaluator tras jest bezstanowy, ocena populacji daje się
zrównoleglić, co umożliwiłoby eksperymenty o~znacznie większym budżecie.

\paragraph{Rozwój metodyki badań.}
Zwiększenie liczby instancji podniosłoby moc testów i~pozwoliło rozstrzygnąć,
czy mała, nieistotna statystycznie przewaga NSGA-II nad NSGA-III jest
systematyczna. Hiperobjętość warto uzupełnić miarami o~odmiennym charakterze
(IGD$^{+}$, wskaźnik $\varepsilon$, miary rozpiętości) oraz zbadać wrażliwość
wyników na położenie punktu referencyjnego. Interesującym kierunkiem jest też
porównanie przy równym czasie procesora zamiast równej liczby ocen oraz analiza
zależności między cechami instancji (np. udziałem komórek zakazanych, pokryciem
radarowym) a~wynikami poszczególnych algorytmów, która mogłaby wyjaśnić dużą
zmienność między instancjami zaobserwowaną w~rozdziale~\ref{ch:analiza}.

\section{Potencjalne obszary zastosowania}
\label{sec:podsumowanie_zastosowania}

Opracowane rozwiązanie ma dwojaki charakter: jest zarazem narzędziem
planistycznym i~środowiskiem badawczym.

Jako narzędzie planistyczne bezpośrednio odpowiada potrzebom misji, w~których
zespół dronów ma dotrzeć do wyznaczonych celów możliwie szybko i~oszczędnie,
unikając przy tym wykrycia. Takie sformułowanie jest właściwe dla zadań
rozpoznawczych i~obserwacyjnych w~zastosowaniach obronnych, gdzie strefy
wykrywalności odpowiadają zasięgom radarów przeciwnika, ale przenosi się także
na zastosowania cywilne: w~akcjach poszukiwawczo-ratowniczych i~monitoringu
infrastruktury strefy ,,wykrywalności'' można interpretować jako obszary
zakłóceń łączności lub ruchu lotniczego, a~w~logistyce dronowej --- jako strefy
o~podwyższonym ryzyku lub ograniczeniach prawnych. Istotną zaletą podejścia
wielokryterialnego jest to, że wynikiem nie jest jedna trasa, lecz zbiór tras
kompromisowych; decyzja o~tym, ile czasu lub energii warto poświęcić dla
zmniejszenia ryzyka, pozostaje po stronie operatora, który dokonuje jej
a~posteriori, mając przed sobą cały front. Należy jednak podkreślić, że model
symulacyjny jest uproszczony (dyskretna siatka, brak dynamiki lotu, planowanie
z~pełną znajomością środowiska), więc zastosowanie praktyczne wymagałoby
integracji z~warstwą sterowania lotem i~walidacji na danych rzeczywistych.

Jako środowisko badawcze symulator wraz z~protokołem eksperymentalnym stanowi
gotowe stanowisko do testowania metaheurystyk dyskretnych: wspólny interfejs
algorytmów, bezstanowy ewaluator, dokładna metryka hiperobjętości oraz skrypty
analizy statystycznej pozwalają włączyć nowy algorytm do porównania bez
ingerencji w~pozostałe komponenty. Pełna odtwarzalność eksperymentów czyni go
również przydatnym w~dydaktyce --- do ilustrowania pojęć dominacji Pareto,
hiperobjętości i~równowagi między eksploracją a~eksploatacją na problemie
o~czytelnej interpretacji fizycznej.
